% Evaluation chapter for REtcd thesis.
%
% Requires the same preamble packages and the Go listings definition as
% design.tex (booktabs, listings, xcolor, hyperref, textcomp). See the header
% comment in design.tex for the recommended \lstdefinelanguage{Go} / \lstset
% block.
%
% Cross-references into the design chapter (e.g. \ref{sec:design-alternatives})
% assume design.tex is \input before this file.

\chapter{Evaluation}
\label{ch:evaluation}

This chapter evaluates REtcd against native etcd along the axis the design chapter identified as REtcd's signature cost: per-operation latency. The evaluation is deliberately scoped to two complementary experiments. Section~\ref{sec:eval-questions} states the questions. Section~\ref{sec:eval-setup} describes the two testbeds. Section~\ref{sec:eval-steady} reports the primary result---a controlled microbenchmark that isolates the apiserver$\rightarrow$storage write path and shows that, in steady state, REtcd's single Redis hop is not only competitive with etcd's in-process bbolt but exhibits a markedly tighter latency tail. Section~\ref{sec:eval-coldstart} reports a preliminary macrobenchmark of Knative cold-start latency under a serverless trace, where the opposite picture emerges in the tail; Section~\ref{sec:eval-coldstart} is explicit about why this result is preliminary and what it can and cannot support. Section~\ref{sec:eval-threats} enumerates the threats to validity.

The headline finding is a two-part one, and the two parts are not in tension: \emph{per isolated operation, REtcd matches or beats etcd; chained across the many sequential writes of a Knative scale-from-zero, REtcd's per-RPC overhead compounds into a tail regression.} The mechanism connecting them---latency amplification across sequential control-plane writes---is the empirical contribution of this chapter.

\section{Evaluation questions}
\label{sec:eval-questions}

The evaluation is organised around three questions:

\begin{itemize}
    \item \textbf{EQ1 --- Steady-state write latency.} For the individual write operations that dominate Kubernetes control-plane traffic (pod create, pod delete), how does REtcd's request latency compare to native etcd, across the full distribution from median to tail?
    \item \textbf{EQ2 --- End-to-end cold-start latency.} Under a realistic Function-as-a-Service workload (Knative Serving driven by a production-derived invocation trace), how does the choice of storage backend affect the latency a client observes, including the scale-from-zero path?
    \item \textbf{EQ3 --- Cost attribution.} Where does any observed difference originate---in the storage substrate's per-operation cost, or in the number and sequencing of operations the control plane issues?
\end{itemize}

EQ1 is answered with a controlled microbenchmark (Section~\ref{sec:eval-steady}) and is the result the thesis stands on. EQ2 is answered with a single preliminary run per backend (Section~\ref{sec:eval-coldstart}); its role is to motivate EQ3 and the cost-attribution argument, not to support a precise quantitative claim. EQ3 is addressed by reading the two results together.

\section{Experimental setup}
\label{sec:eval-setup}

Two testbeds are used, chosen to isolate different parts of the system.

\subsection{Microbenchmark: pod-churn against a kwok-backed control plane}
\label{sec:eval-micro-setup}

The microbenchmark (\texttt{load-test/load\_test.py}) is an open-loop Poisson load generator that drives pod-lifecycle operations directly against the kube-apiserver. Each \emph{lifecycle} is one pod \texttt{CREATE} followed, after a fixed hold time, by a \texttt{DELETE}; new lifecycles arrive as a Poisson process with mean rate $\lambda$. The generator records the wall-clock latency of every apiserver request to CSV and computes per-trial validity (a trial is invalid if the non-conflict failure rate exceeds 1\%).

The critical methodological choice is that pods are scheduled onto \textbf{kwok} virtual nodes (\texttt{nodeSelector: type=kwok}), not real kubelets. kwok acknowledges pod binding without starting containers, which removes scheduling, image-pull, and kubelet latency from the measurement. What remains on the measured path is exactly the apiserver$\rightarrow$storage write: admission, validation, and the etcd \texttt{Txn}/\texttt{Range}/\texttt{DeleteRange} that the backend services. This isolates EQ1's variable---the storage backend---from the rest of the control plane.

\paragraph{Parameters.}
Target rate $\lambda = 30$ lifecycles/s; 5\,s unmeasured warmup followed by a 30\,s measured window; pod hold time 1\,s; 3 independent trials per backend. After the measured window the generator drains in-flight lifecycles so that deletes are recorded. These settings yield approximately 2{,}600 create and 2{,}600 delete observations per backend in the measured window.

\paragraph{Backends compared.}
Two single-instance configurations, identical in every respect except the storage backend:

\begin{enumerate}
    \item \textbf{etcd} --- a stock single-node etcd, the kind cluster default (\texttt{kind-etcd-baseline}).
    \item \textbf{REtcd} --- REtcd backed by a single local Redis with AOF persistence (\texttt{everysec}).
\end{enumerate}

Both are single-instance: the etcd baseline runs no Raft quorum, so this is a \emph{fair single-instance comparison}, not a comparison against highly-available etcd. The comparison reported in Section~\ref{sec:eval-steady} uses, for each backend, a run in which all three trials were valid with a zero non-conflict failure rate; earlier REtcd runs that contained an invalid trial (driven into overload, peak in-flight 313) or a configuration fault are excluded and noted in Section~\ref{sec:eval-threats}.

\begin{figure}[h]
    \centering
    % TODO: Replace placeholder with the microbenchmark topology diagram.
    % Boxes: load_test.py (open-loop Poisson generator) -> kube-apiserver
    %        -> {etcd | REtcd -> Redis}
    % Separate box "kwok virtual nodes" attached to the scheduler, annotated
    %   "binds pods without running containers --- removes kubelet latency from
    %    the measured path".
    % Annotate generator->apiserver edge: "measured: per-request latency".
    % TODO: state the host hardware (CPU model, cores, RAM) and the
    %       kind / Kubernetes / etcd / Redis versions used.
    \fbox{\parbox{0.9\linewidth}{\centering\vspace{2em}\textit{[Placeholder: microbenchmark topology --- see TODO in source]}\vspace{2em}}}
    \caption{Microbenchmark topology. kwok virtual nodes bind pods without running containers, isolating the apiserver$\rightarrow$storage write path.}
    \label{fig:eval-micro-topo}
\end{figure}

\subsection{Macrobenchmark: Knative cold-start on CloudLab}
\label{sec:eval-macro-setup}

The macrobenchmark measures end-to-end client-observed response time for serverless function invocations, including the scale-from-zero path that dominates FaaS cold-start latency. It runs on a 3-node CloudLab cluster (hardware type \textbf{xl170}, Ubuntu 22.04) running Kubernetes 1.32.13, Knative Serving 1.15 with the Kourier ingress, and is driven by the \texttt{invitro} loader replaying the \texttt{azure\_10} invocation trace over a 5-minute window. The same cluster, trace, and window are used for both backends; only the storage backend is swapped.

This testbed answers EQ2 but, as Section~\ref{sec:eval-coldstart} makes explicit, the run reported here is a single smoke run per backend with a small number of in-window invocations, and is presented as a preliminary observation rather than a characterised distribution.

\section{Steady-state write latency}
\label{sec:eval-steady}

Table~\ref{tab:eval-steady} reports the latency distribution for pod create and delete under the microbenchmark, pooled across the three valid trials of each backend. Figure~\ref{fig:eval-cdf} shows the corresponding latency CDFs; Figure~\ref{fig:eval-bars} shows per-percentile bars with error bars across trials.

\begin{table}[h]
    \centering
    \caption{Apiserver request latency (ms) under kwok-isolated pod churn at $\lambda = 30$/s, pooled over three valid trials per backend. Lower is better.}
    \label{tab:eval-steady}
    \begin{tabular}{@{}llrrrrrrr@{}}
        \toprule
        \textbf{Backend} & \textbf{Op} & \textbf{n} & \textbf{p50} & \textbf{p90} & \textbf{p99} & \textbf{p99.9} & \textbf{mean} & \textbf{max} \\
        \midrule
        etcd  & create & 2629 & 6.15 &  8.71 & 12.94 & 41.15 & 6.60 &  92.14 \\
        REtcd & create & 2588 & 5.66 &  8.07 & 10.54 & 16.82 & 5.93 &  24.78 \\
        \addlinespace
        etcd  & delete & 2617 & 7.22 & 10.30 & 14.23 & 62.42 & 7.75 & 130.53 \\
        REtcd & delete & 2597 & 7.88 & 10.45 & 13.34 & 23.53 & 8.07 &  32.13 \\
        \bottomrule
    \end{tabular}
\end{table}

\begin{figure}[h]
    \centering
    % TODO: Replace with latency CDF (log-x), one panel for create, one for delete.
    % Source: load-test/plots/cdf_create.pdf and load-test/plots/cdf_delete.pdf
    % Annotation: "REtcd body tracks etcd; REtcd tail is bounded well below etcd's."
    \fbox{\parbox{0.9\linewidth}{\centering\vspace{2em}\textit{[Placeholder: latency CDFs (create, delete) --- see TODO in source]}\vspace{2em}}}
    \caption{Latency CDF (log-$x$) for pod create and delete, etcd vs REtcd. The REtcd body tracks etcd; the REtcd tail is bounded well below etcd's.}
    \label{fig:eval-cdf}
\end{figure}

\begin{figure}[h]
    \centering
    % TODO: Replace with per-percentile bar chart (p50/p90/p99/p99.9),
    % mean +/- std across the three trials, etcd vs REtcd, for create and delete.
    % Source: load-test/plots/pXX_bar_create.pdf and load-test/plots/pXX_bar_delete.pdf
    \fbox{\parbox{0.9\linewidth}{\centering\vspace{2em}\textit{[Placeholder: percentile bar chart --- see TODO in source]}\vspace{2em}}}
    \caption{Latency percentiles (p50/p90/p99/p99.9), mean $\pm$ standard deviation across the three trials, etcd vs REtcd.}
    \label{fig:eval-bars}
\end{figure}

Three observations follow from Table~\ref{tab:eval-steady}.

\paragraph{The single Redis hop does not penalise the median.}
On create, REtcd's median is 5.66\,ms against etcd's 6.15\,ms---REtcd is in fact 8\% \emph{faster} at p50. On delete the median favours etcd slightly (7.22\,ms vs 7.88\,ms, a 9\% difference). The design chapter predicted a per-RPC kernel-IPC cost of order 50--100\,\textmu{}s for the Redis hop over TCP localhost; against a median request latency of several milliseconds, that cost is below the noise floor of this experiment, and the two backends are within $\pm$10\% of each other at the median for both operations. The extra network hop, in steady state, is not visible at the median.

\paragraph{REtcd's tail is substantially tighter.}
The divergence between the backends is in the tail, and it runs in REtcd's favour. At p99.9, REtcd's create latency is 16.8\,ms against etcd's 41.1\,ms (59\% lower), and its delete latency is 23.5\,ms against etcd's 62.4\,ms (62\% lower). The worst observed latency tells the same story more starkly: REtcd's maximum create/delete latencies (24.8\,ms / 32.1\,ms) are roughly a quarter of etcd's (92.1\,ms / 130.5\,ms). REtcd's tail is not merely competitive---it is bounded well below etcd's across both operations.

\paragraph{A plausible mechanism for the tail difference.}
This result is consistent with the two backends' persistence architectures, though the present experiment does not instrument the cause directly. Native etcd commits through a write-ahead log with periodic \texttt{fsync}, a bbolt B+-tree whose page rebalancing and freelist management occur in bursts, and a background MVCC compaction cycle; each of these can inject occasional multi-tens-of-milliseconds stalls into individual requests, producing the heavy tail seen in etcd's p99.9 and max. Redis under \texttt{appendonly everysec} defers durability to a once-per-second background flush off the request path, so individual writes see comparatively uniform service time. Attributing the tail difference to these mechanisms specifically would require per-request server-side tracing on both backends and is left to future work; the claim made here is the measured one---REtcd's tail distribution is tighter---not a claim about its cause.

Taken together, these answer EQ1: for the isolated write operations that constitute the bulk of control-plane traffic, REtcd is competitive with native etcd at the median and superior in the tail, on a fair single-instance comparison.

\section{Cold-start latency (preliminary)}
\label{sec:eval-coldstart}

The microbenchmark isolates a single operation. A Knative scale-from-zero does the opposite: it chains a sequence of dependent control-plane writes---autoscaler decision, SKS update, Deployment scale, scheduler binding, kubelet status, EndpointSlice publication---each of which is itself one or more apiserver writes that cross apiserver$\rightarrow$backend. Table~\ref{tab:eval-coldstart} reports a single preliminary run per backend on the Section~\ref{sec:eval-macro-setup} macrobenchmark.

\begin{table}[h]
    \centering
    \caption{Preliminary end-to-end invocation response time (ms) on the 3-node CloudLab cluster, \texttt{azure\_10} trace, 5-minute window. \textbf{This is a single smoke run per backend; see caveats below.}}
    \label{tab:eval-coldstart}
    \begin{tabular}{@{}lrrrrrrr@{}}
        \toprule
        \textbf{Backend} & \textbf{n} & \textbf{min} & \textbf{median} & \textbf{mean} & \textbf{p95} & \textbf{p99} & \textbf{max} \\
        \midrule
        etcd  & 43 & 5.0 & 84.8 &  1358.8 &   8148 &  10042 &  10042 \\
        REtcd & 44 & 5.5 & 73.9 & 26125.8 & 107500 & 158639 & 158639 \\
        \bottomrule
    \end{tabular}
\end{table}

\begin{figure}[h]
    \centering
    % TODO: Replace with response-time CDF, etcd vs REtcd, on the macrobenchmark.
    % Source: results-retcd-3node-smoke/smoke_responseTime_cdf.png
    % Annotation: "bodies coincide to ~0.6; REtcd tail diverges, dominated by
    %              grpcConnEstablish during scale-from-zero."
    \fbox{\parbox{0.9\linewidth}{\centering\vspace{2em}\textit{[Placeholder: response-time CDF --- see TODO in source]}\vspace{2em}}}
    \caption{Response-time CDF, etcd vs REtcd, on the macrobenchmark. The bodies coincide to roughly the 60th percentile; the REtcd tail diverges, dominated by \texttt{grpcConnEstablish} during scale-from-zero.}
    \label{fig:eval-coldstart-cdf}
\end{figure}

Read with appropriate caution, two things are visible. First, the \textbf{body} of the distribution is indistinguishable between backends: up to roughly the 60th percentile the two CDFs coincide, and REtcd's median (73.9\,ms) is actually below etcd's (84.8\,ms)---consistent with the steady-state finding of Section~\ref{sec:eval-steady}, since warm invocations exercise the same single-operation path. Second, the \textbf{tail} diverges sharply: REtcd's p99 is 158.6\,s against etcd's 10.0\,s. Inspection of the slow REtcd invocations attributes the time to \texttt{grpcConnEstablish} during scale-from-zero---the function executes normally once its pod is running; the wait is for the scaling pipeline to bring a scaled-to-zero service back up. Because that pipeline is a chain of $\sim$7--8 sequential apiserver writes, each crossing apiserver$\rightarrow$REtcd$\rightarrow$Redis, a per-RPC overhead that is invisible in isolation (Section~\ref{sec:eval-steady}) is multiplied by the chain length and surfaces in the cold-start tail. This is the latency-amplification mechanism named at the start of the chapter, and it is the bridge to EQ3: the cost is not the substrate's per-operation latency in isolation but the interaction of that latency with the sequential, write-heavy structure of the cold-start path.

\paragraph{Why this result is preliminary.}
Four limitations bar any precise quantitative claim from Table~\ref{tab:eval-coldstart}, and they are stated plainly:

\begin{enumerate}
    \item \textbf{Sample size.} With $n \approx 44$ in-window invocations per backend, the reported p99 is effectively the single worst observation; it carries no confidence interval. A characterised tail requires a larger trace (\texttt{azure\_50} or larger) over a 15--30 minute window.
    \item \textbf{One run per backend.} There is no across-run variance characterisation, so the difference between backends cannot be separated from run-to-run noise. At least three runs per backend are required before stating a percentage tail regression.
    \item \textbf{Scale.} Three nodes is not the reference testbed; the closest comparable work (KUBEDIRECT, Section~\ref{sec:design-alternatives}) evaluates at 80 nodes. Triangulating against published numbers requires matching that scale.
    \item \textbf{What the tail measures.} The 158\,s figure is dominated by the autoscaler control loop during scale-from-zero, not by storage service time per se; it is best read as evidence for the \emph{amplification mechanism}, not as a measurement of REtcd's storage latency.
\end{enumerate}

Accordingly, Section~\ref{sec:eval-coldstart} is presented as motivating evidence for EQ3 and as a pointer to the future work, not as a headline performance number.

\section{Threats to validity}
\label{sec:eval-threats}

\paragraph{Internal validity.}
The microbenchmark's isolation rests on kwok faithfully removing kubelet-side latency while preserving the apiserver$\rightarrow$storage path; if kwok's admission shortcuts alter the apiserver's write pattern relative to a real kubelet, the isolated path would not be representative. The excluded REtcd runs (one trial driven into overload at peak in-flight 313; one run with a 50\% failure rate attributable to a configuration fault) indicate that the generator can enter open-loop backlog under stress; the reported run was checked to have bounded in-flight concurrency (peak $\approx$ 47, comparable to etcd's $\approx$ 48), so the two reported backends were measured under comparable load, not with REtcd advantaged by a lighter offered load.

\paragraph{Construct validity.}
EQ1 measures apiserver request latency, which includes admission and validation common to both backends; it is therefore a conservative estimate of the \emph{relative} backend difference, since the shared overhead dilutes it. The Section~\ref{sec:eval-steady} tail result is reported as a measured distribution, not a causal claim about bbolt vs AOF; the mechanism in Section~\ref{sec:eval-steady} is a hypothesis pending server-side tracing.

\paragraph{External validity.}
Both experiments are single-instance; nothing here speaks to highly-available, Raft-quorate etcd, which pays inter-node consensus RTT that REtcd's single-instance design does not attempt to match. The microbenchmark host hardware (TODO, Section~\ref{sec:eval-micro-setup}) and the 3-node macrobenchmark scale both differ from the reference FaaS testbed; generalising the cold-start result to production scale is explicitly out of reach of the present data.

\paragraph{Statistical validity.}
The Section~\ref{sec:eval-steady} result has three trials per backend and $\sim$2{,}600 observations per cell, sufficient to read medians and to compare tails up to p99.9 with the per-trial error bars of Figure~\ref{fig:eval-bars}. The Section~\ref{sec:eval-coldstart} result, with $n \approx 44$ and one run per backend, supports no statistical claim and is labelled preliminary throughout.

These threats define the shape of the remaining work: enlarging the macrobenchmark to a characterised distribution (more invocations, multiple runs, larger scale), and instrumenting the server side to confirm the tail-latency mechanism hypothesised in Section~\ref{sec:eval-steady}.
